% Generated draft: requires mathematical and source review.
\documentclass[11pt]{article}
\usepackage[margin=25mm]{geometry}
\usepackage{fontspec}
\setmainfont{Latin Modern Roman}
\usepackage{amsmath,amssymb,mathtools,unicode-math}
\setmathfont{Latin Modern Math}
\usepackage{enumitem,needspace,xurl,hyperref}
\setlength{\parindent}{0pt}
\setlength{\parskip}{5pt}
\setlength{\emergencystretch}{3em}
\newcommand{\N}{\mathbb{N}}
\newcommand{\Z}{\mathbb{Z}}
\newcommand{\Q}{\mathbb{Q}}
\newcommand{\R}{\mathbb{R}}
\newcommand{\C}{\mathbb{C}}
\newcommand{\F}{\mathbb{F}}
\newcommand{\A}{\mathbb{A}}
\newcommand{\PP}{\mathbb{P}}
\newcommand{\E}{\mathbb{E}}
\newcommand{\eps}{\varepsilon}
\newcommand{\dd}{\,\mathrm{d}}
\begin{document}
\section*{UnsolvedMath ID: 800642}
\textbf{Source problem code:} catalogue entry
\subsection{Definitions and mathematical statement}
The source does not specify the ambient class of spaces or its conventions for covers. The following gives a formulation using standard countable-cover conventions; agreement with the cited source requires review.

For a topological space \(X\), a cover is a countable family \(\mathcal U\) of subsets of \(X\) such that \(\bigcup\mathcal U=X\) and \(X\notin\mathcal U\). It is a \(\gamma\)-cover if it is infinite and each point of \(X\) belongs to all but finitely many members. It is a \(\tau\)-cover if each point belongs to infinitely many members and, for every \(x,y\in X\), at least one of the sets
\[
\{U\in\mathcal U:x\in U,\ y\notin U\},\qquad
\{U\in\mathcal U:y\in U,\ x\notin U\}
\]
is finite. Write \(\Gamma(X)\) and \(\mathrm T(X)\) for the families of open \(\gamma\)-covers and open \(\tau\)-covers, respectively. Write \(\mathcal B_\Gamma(X)\) and \(\mathcal B_{\mathrm T}(X)\) for the corresponding families whose members are Borel subsets of \(X\).

For cover families \(\mathcal A(X)\) and \(\mathcal B(X)\), define the following properties:
\begin{itemize}
\item \(\mathsf S_1(\mathcal A,\mathcal B)\): for every sequence \((\mathcal U_n)_{n\in\N}\) in \(\mathcal A(X)\), there are choices \(U_n\in\mathcal U_n\) such that \(\{U_n:n\in\N\}\in\mathcal B(X)\).
\item \(\mathsf S_{\mathrm{fin}}(\mathcal A,\mathcal B)\): for every sequence \((\mathcal U_n)_{n\in\N}\) in \(\mathcal A(X)\), there are finite subfamilies \(\mathcal F_n\subseteq\mathcal U_n\) such that \(\bigcup_{n\in\N}\mathcal F_n\in\mathcal B(X)\).
\item \(\mathsf U_{\mathrm{fin}}(\mathcal A,\mathcal B)\): for every sequence \((\mathcal U_n)_{n\in\N}\) in \(\mathcal A(X)\) such that no \(\mathcal U_n\) has a finite subcover, there are finite subfamilies \(\mathcal F_n\subseteq\mathcal U_n\) such that \(\{\bigcup\mathcal F_n:n\in\N\}\in\mathcal B(X)\).
\end{itemize}
Here \(\N=\{1,2,\ldots\}\), and each displayed output family is treated as a set of subsets, without counting repetitions.

\textbf{Question (Tsaban's list, Problem 15; context \textnormal{[AddQuad]}).} For each of the properties
\[
\mathsf S_1(\mathrm T,\mathrm T),\quad
\mathsf S_{\mathrm{fin}}(\mathrm T,\mathrm T),\quad
\mathsf S_1(\Gamma,\mathrm T),\quad
\mathsf S_{\mathrm{fin}}(\Gamma,\mathrm T),\quad
\mathsf U_{\mathrm{fin}}(\Gamma,\mathrm T),
\]
and each corresponding Borel version obtained by replacing \(\Gamma\) and \(\mathrm T\) by \(\mathcal B_\Gamma\) and \(\mathcal B_{\mathrm T}\), is that property preserved under finite unions?

Explicitly, for each such property \(P\), ask whether, for every positive integer \(m\) and every union \(X=\bigcup_{i=1}^{m}X_i\) allowed by the intended domain, with each \(X_i\) carrying the subspace topology from \(X\),
\[
\bigl(\forall i\in\{1,\ldots,m\}\;X_i\text{ satisfies }P\bigr)
\quad\Longrightarrow\quad X\text{ satisfies }P.
\]
The intended domain of these unions is not specified in the supplied record. In particular, whether the question concerns sets of reals, a specified class of spaces, or arbitrary topological spaces remains unresolved; the displayed implication is a formulation template pending that specification.
\subsection{Short English statement}
For each of five listed selection principles, and for its Borel version, ask whether every finite union of spaces having the property also has the property. The intended class of spaces and the source's exact cover conventions require confirmation. The registry labels this problem open; that label is supplied metadata, not an independently verified status assessment.
\subsection{Sources}
\begin{itemize}
\item Local UnsolvedMath registry, permanent ID \(800642\), exact source problem code \texttt{catalogue entry}; title: \emph{Selection Principles Milestone Problems --- Problem 15 ([AddQuad])}. The supplied record attributes the problem to Tsaban, with proposed year 2003, and labels its status \texttt{open}.
\item Supplied source reference: Tsaban, \emph{Selection Principles in Mathematics: A Milestone of Open Problems} (2003), Problem 15. Supplied URL: \url{https://arxiv.org/abs/math/0312182}. The record reports access on 2026-10-16 and extraction method \texttt{arxiv\_latex}; these are supplied provenance statements, not a claim of independent access or verification.
\item Supplied attribution/context: \texttt{[AddQuad]}. The record does not expand this reference key or supply its bibliographic details.
\end{itemize}
\end{document}
